% =============================================================================
% main.tex – Selection Panel Firmware: Technische Dokumentation
% =============================================================================
% Architektur-Entscheidung: main.tex bleibt ein reiner Orchestrator.
% Inhalt liegt in content/*.tex → Single Responsibility, einfachere Wartung.
% =============================================================================
\documentclass{techdoc}
% Optionen für Ausgabevarianten (Kompilierung identisch, Rendering unterschiedlich):
% \documentclass[print]{techdoc}  % Links unsichtbar → sauberer Ausdruck
% \documentclass[draft]{techdoc}  % Overfull-Boxen + TODOs sichtbar → Frühwarnsystem

% =============================================================================
% METADATEN
% =============================================================================
% SSOT-Prinzip: Alle Metadaten zentral, damit Titelseite, Kopfzeilen und
% PDF-Properties automatisch konsistent bleiben (keine manuelle Synchronisation).
% -----------------------------------------------------------------------------
\title{Interaktives Auswahlpanel}
\subtitle{10× Taster/LED}
\tagline{\glqq Drücken -- Leuchten\grqq}
\scope{ESP32-S3 \textbullet\ Raspberry Pi 5 \textbullet\ Web-Dashboard}
\stack{CD4021BE \textbullet\ 74HC595 \textbullet\ PlatformIO \textbullet\ Python/aiohttp \textbullet\ WebSocket \textbullet\ HTML/JS \textbullet\ systemd}
\version{2.5.3}
\date{20.01.2026}
\author{Jan Unger}
\credits{Claude Opus 4.5, ChatGPT 5.2}
\coverimage{cover-embedded-projekt.pdf}

% Bilder liegen in ./images/ und ./covers/
\graphicspath{{images/}{covers/}}

% Kurzform für Kopfzeile, da vollständiger Titel zu lang wäre.
\setheadertitle{Interaktives Auswahlpanel}

% PDF-Metadaten für Suchmaschinen und Dokumentenverwaltung.
\hypersetup{
  pdftitle={Interaktives Auswahlpanel – 10× Taster/LED},
  pdfauthor={Jan Unger},
  pdfsubject={Embedded-Projekt mit ESP32-S3 und Schieberegistern},
  pdfkeywords={ESP32-S3, Raspberry Pi 5, FreeRTOS, WebSocket, React}
}

% Externes PDF einbinden → Cover wurde in separatem Tool gestaltet (Designtrennung).
\usepackage{pdfpages}

% =============================================================================
% DOKUMENT
% =============================================================================
\begin{document}

% Titelseite generieren (nutzt alle oben definierten Metadaten)
\maketitlepage

% Römische Ziffern im Vorspann → Konvention für Verzeichnisse, damit
% Seitenzahl 1 erst beim eigentlichen Inhalt beginnt.
\pagenumbering{roman}
\tableofcontents
\clearpage
\pagenumbering{arabic}

% =============================================================================
% TEIL I: THEORETISCHE GRUNDLAGEN
% =============================================================================
% Gruppierung mit {} → lokaler Scope, damit Formatierungen (z.B. \LARGE)
% nicht in nachfolgende Abschnitte durchsickern.
{%
  \clearpage
  \vspace*{2cm}
  \begin{center}
    {\sffamily\bfseries\LARGE Teil I}\\[0.5em]
    {\sffamily\bfseries\Huge Theoretische Grundlagen}\\[1em]
    {\large\color{MutedText} Mikrocontroller, Mindset, Entwicklungsprozess}
  \end{center}
  \vspace{1cm}
  % Manueller TOC-Eintrag, da Teil-Seite kein \section ist.
  \addcontentsline{toc}{section}{Teil I: Theoretische Grundlagen}
}

\begin{infobox}[Ziel]
Dieses Kapitel vermittelt die fundamentalen Konzepte der Embedded-Entwicklung:
Mikrocontroller-Architektur, ressourcenbewusstes Denken und professionelle
Entwicklungspraktiken.
\end{infobox}

\clearpage

% Kapitel 1-3 bauen didaktisch aufeinander auf:
% Hardware-Verständnis → Denkweise → Werkzeuge
\input{content/embedded-grundlagen}
\clearpage

\input{content/embedded_mindset}
\clearpage

\input{content/embedded-entwicklung}
\clearpage

% =============================================================================
% TEIL II: DAS SELECTION-PANEL-PROJEKT
% =============================================================================
% Trennung Theorie/Praxis → Leser können Teil I überspringen, wenn
% Grundlagen bekannt sind.
{%
  \clearpage
  \vspace*{2cm}
  \begin{center}
    {\sffamily\bfseries\LARGE Teil II}\\[0.5em]
    {\sffamily\bfseries\Huge Das Selection-Panel-Projekt}\\[1em]
    {\large\color{MutedText} Architektur, Protokolle, Web-Integration}
  \end{center}
  \vspace{1cm}
  \addcontentsline{toc}{section}{Teil II: Das Selection-Panel-Projekt}
}

\begin{infobox}[Projektüberblick]
Dieser Teil dokumentiert das Selection-Panel-Projekt: Systemarchitektur,
Kommunikationsprotokolle und Web-Integration. Die Kapitel bauen aufeinander auf
-- vom Gesamtbild zu den technischen Details.
\end{infobox}

\clearpage

% Kapitel 4-6: Top-Down-Struktur (Architektur → Protokoll → UI)
% ermöglicht schrittweises Einarbeiten ohne Vorwärtsreferenzen.
\input{content/projekt}
\clearpage

\input{content/Serial-Protokoll}
\clearpage

\input{content/Frontend-Backend}
\clearpage

% =============================================================================
% ANHANG
% =============================================================================
% \appendix setzt Zähler auf A, B, C → klare Trennung Hauptteil/Referenzmaterial.
\appendix

{%
  \clearpage
  \vspace*{2cm}
  \begin{center}
    {\sffamily\bfseries\Huge Anhang}\\[1em]
    {\large\color{MutedText} Glossar und Referenzen}
  \end{center}
  \vspace{1cm}
  \addcontentsline{toc}{section}{Anhang}
}

\clearpage

% Glossar am Ende → Leser können bei Bedarf nachschlagen, ohne den
% Lesefluss im Haupttext zu unterbrechen.
\input{content/embedded-glossar}

\end{document}
